\documentclass[10pt]{article}

\usepackage[utf8]{inputenc}

\usepackage[T1]{fontenc}

\usepackage{amsmath}

\usepackage{amsfonts}

\usepackage{amssymb}

\usepackage[version=4]{mhchem}

\usepackage{stmaryrd}

\usepackage{bbold}



\begin{document}

Значительно больше, чем в общем случае, известно о коммутативных алгебраических У. г. (см. [4]). Если $\operatorname{char} K=0$, то әто в точности алгебраич. группы, изоморфные $\mathbf{G}_{a} \times \ldots \times \mathbf{G}_{a}$; при этом изоморфизм $\mathbf{G}_{a} \times \ldots$ $\times \mathbf{G}_{a} \rightarrow U$ задается әкспоненциальным отображением. Если же char $K=p>0$, то связные коммутативные алгебраические У. г. $U$ - әто в точности связные коммутативные алгебраич. групшы, являющиеся $p$-группами. В этом случае $U$ не обязательно изоморфна $\mathbf{G}_{a} \times \ldots$ $\times \mathbf{G}_{a}$ : для әтого необходимо и достаточно, чтобы $g p=e$ для любого $g \in U$. В общем же случае $U$ изогенна произведению нек-рых специальных групі (т. н. групіп Витта, см. [2]).



Если $H$ и $U$ - связные алгебраические У . г. и $H \subset U$, то многообразие $U / H$ изоморфно аффинному пространству. Любая орбита алгебраической У. г. автоморфизмов афффинного алгебраич. многообразия $X$ замкнута в $X[5]$.



Лит. - [1] Б о р е л ь А., Линейные алгебраические группы, пер. с англ., М., 1972, [2] С е р р Ж.. П., Алгебраические группы и поля классов, пер с франц., М., 1968; [3] Х а м ф р и Д ж. Линейные алгебраические группы, пер. с англ., M., 1980, [4] Unipotent algebraic groups, B.- [a. o.], 1974, [5] S t e i nb e r g R., Conjugacy classes in algebraic groups,

1974. УНИПОТЕНТНАЯ МАТРИЦА - квадратная матрица $A$ над кольцом, для к-рой матрица $A-E_{n}$, где $n-$ порядок матрицы $A$, нильпотентна, то есть $\left(A-E_{n}\right)^{n}=$ $=0$. Матрица над полем порядка $n$ унипотентна тогда и только тогда, когда ее характеристический многочлен есть $(x-1)^{n}$.



Группа матриц наз. у н и п о т е н т н о й, если каждая ее матрица унипотентна. Любая унипотентная подгрупша в $G L(n, F)$, где $F$ - поле, сопряжена в $G L(n, F)$ с нек-рой подгруппой специальной треугольной группы (т е о р е м а $К$ о л ч и н а). Это утверждение справедливо и для унипотентных групп над телом, если характеристика тела либо равна 0 , либо больше некрого $\gamma(n)$. Д. А. Супруненко.



УНИПОТЕНТНЫЙ ЭЛЕМЕНТ - элемент $g$ линейной алгебраич. группы $G$, совпадающий с унипотентной компонентой $g_{u}$ своего Жордана разложения в группс $G$. Если реализовать $G$ как замкнутую подгруппу группы $G L(V)$ автоморфизмов конечномерного векторного пространства $V$ над основным алгебраически замкнутым полөм $K$, то У. э. $g$ - это в точности такие әлементы, что $(1-g)^{n}=0, n=\operatorname{dim} V$, или, что их матрицы в нек-ром базисе пространства $V$ являются верхне-треугольными с единицами на диагонали. Множество $U(G)$ всех У. ә. в $G$ замкнуто в топологии Зариского, а если $G$ определена над подполем $k \subset K$, то и $U(G)$ определено над $k$. Если char $K=0$, то всякий У. ә. $g$ имеет бесконечный порядок. В этом случае наименьшая алгебраич. подгруппа в $G$, содержащая $g$, является одномерной унипотентной әруппой. Если же char $K=p>0$, то $g$ будет У. ә. в точности тогда, когда он имеет конечный порядок, равный $p^{t}$ для нек-рого $t \geqslant 0$. Связная группа не содержит $У$. э. $g \neq e$ тогда и только тогда, когда она является алгебрачческим тором.



В терминах У. э. может быть дан критерий анизотропности (см. Анизотропнал группа).



У. ә. играют важную роль в теории дискретных подepynn алгебраич. групп и групп Ли. Наличие У. э. в дискретных группах $\Gamma$ движений симметрических пространств, имеющих некомпактную фундаментальную область конечного объема, является важным средством изучения структуры таких групп и их канонических фундаментальных областей (см. [5]); существование У. ә. в таких $\Gamma$ доказано в [4].



Многообразие $U(G)$ инвариантно относительно внутренних автоморфизмов группы $G$. Пусть $G$ связна и полупроста. Тогда число классов согряженных У. ә. конечно и для каждой простой $G$ известно полное их описание (а также описание централизаторов У. э.), см. [7]. В классич. группах такое описание получается с использованием жордановой формы матрицы [2]. Напр., для группы $G=S L_{n}(K)$ существует биекция между классами сопряженных У. э. и разбиениями $\left(m_{1}, \ldots\right.$, $m_{s}$ ) числа $n$ в сумму положительных целых слагаемых $m_{i}, m_{1} \geqslant \ldots \geqslant m_{s}$. Если $\lambda=\left(m_{1}, . ., m_{s}\right)$ и $\mu=\left(l_{1}, \ldots\right.$, $l_{t}$ ) - два разбиения числа $n$, то класс, отвечающий $\dot{\lambda}$, содержит в своем замыкании класс, отвечающий $\mu$, в точности тогда, когда $\sum_{i=1}^{j} m_{i} \geqslant \sum_{i=1}^{j} l_{i}$ для всех $j$. Размерность класса, отвечающего разбиению $\left(m_{1}, \ldots\right.$, $m_{s}$ ) (как алгебраического многообразия), равна $n^{2}-\sum_{i j} \min \left(m_{i}, m_{j}\right)$.



Множество всех простых точек алгебраического многообразия $U(G)$ образует один класс сопряженных У. ә.-рег у ля рны е У. э. Если $G$ проста, то многообразие особых точек многообразия $U(G)$ также содержит открытый в топологии Зариского класс сопряженных У. э.- с у б р е г у л я р ны е У. э. По поводу изучения особых точек многообразия $U(G)$ см. также [6].



JIит [1] Б о р е л ь А , Линейные алгебраические группы, пер. с англ., М., 1972; [2] Семинар по алгебраическим группам, пер. с англ., М., 1973, [3] Х а м ф р и Д ж., Јинейные алгебраические группы, пер. с англ., М., 1980, [4] К а ж д а н Д. А , [5] С е л в б е р г А., "Математика", 1972, т. 16, в. 4, с. $72-89$; [6] S l o d o w y P., Simple singularities and simple algebraic groups, B.- [a. o ], 1980; [7] S p a l t e n s t e i n N., Classes unipotentes et sousgroupes de Borel, B.- [a. o ], 1982 .



D. B. JI. Попов. УНИРАЦИОНАЛЬНОЕ МНОГООБРАЗИЕ - аЛгебраическое многообразие $X$ над полем $k$, для к-рого существует такое рациональное отображение проективного пространства $\varphi: P^{n} \rightarrow X$, что $\varphi_{(}\left(P^{n}\right)$ плотно в $X$ и расширение полей рациональных функций $k\left(P^{n}\right) /$ $k(X)$ сепарабельно. Д ругими словами, $k(X)$ имеет сепарабельное чисто трансцендентное расширение.



У. м. близки к рациональныл многообразиял. Напр., на У. м. нет регулярных дифференциальных форм, $H^{0}\left(X, \Omega_{X}^{p}\right)=0$ при $p \geqslant 1$. Вопрос о совпадении понятий рационального и унирационального многообразия наз. Люрота проблемой.



Лит. [1] III а ф а р е в и ч И. Р., Основы алгебраической геометрии, М., 1972. Вик. С. Куликов.



УНИТАРНАЯ ГРУППА о т н о с и т е Ј ь н о ф о рм ы $f$ - группа $U_{n}(K, f)$ всех линейных преобразований $\varphi \quad n$-мерного правого линейного пространства $V$ над телом $K$, сохраняющих фиксированную невырожденную полуторалинейную (относительно инволюции $J$ тела $K$ ) форму $f$ на $V$, т. е. таких $\varphi$, что



$f(\varphi(v), \varphi(u))=f(v, u) \quad \forall v, u \in V$.



У. г. принадлежит к числу классических арупn. Частными случаями У. г. являются симплектическая группа (в этом случае $K$ - поле, $J=1$ и $f$ - знакопеременная билинейная форма) и ортогональная группа ( $K$ - поле, $\operatorname{char} K \neq 2, J=1, f$ - симметрическая билинейная форма). Далее, пусть $J \neq 1$ и $f$ обладает свойством $(T)$ (см. Bumma теорема). Умножая $f$ на подходящий скаляр, можно, не меняя У. г., добиться того, чтобы $f$ стала эрмитовой формой, а меняя, сверх того, $J,-$ чтобы $f$ стала косоәрмитовой формой.



Если исключить случай $n=2, \quad K=\mathbb{F}_{4}$, то всякий элемент У.г. $U_{n}(K, f)$ является произведением не более чем $n+1$ к в а з и о т р а ж е н и й (т. е. преобразований, оставляющих на месте все элементы какой-либо неизотропной гиперплоскости в $V$ ). Центр $Z_{n}$ У. г. $U_{n}(K, f)$ состоит из всех гомотетий пространства $V$ вида $x \mapsto x \gamma, \gamma \in K, \gamma^{J} \gamma=1$.



Пусть $v-$ индекс Витта фороы $f$. Если $v \neq 0$, то удобно считать $f$ косоэрмитовой. Пусть $T_{n}(K, f)-$ нормальный делитель в $U_{n}(K, f)$, порожденный у н и т а р-





\end{document}